\chapter{Synthetic biology and metabolic engineering}
\label{vol06:ch09}

\epigraph{Engineering biology will become a serious discipline only
when it becomes possible to design and reliably construct
biological systems.}{Drew Endy, in the founding paper for the
engineering programme that became modern synthetic biology
\cite[p.~449]{paper:v6c09-endy-foundations-2005}.}

\chapmeta{Half-life: medium-life. The platforms, tools, and tooling
turnover in this field is on the order of five years; the underlying
biology is durable. Archetypes: control, ethics. Exercise target: 25. Project track: simulation and paper design. Hazard class: medium (no wet-lab biosafety required for the chapter projects, but the field's own work runs to BSL-2 and BSL-3 routinely). Reader path default: standard, with sections explicitly tagged. Named cases: He Jiankui 2018 CRISPR-baby case (registry entry \texttt{acc:he-jiankui-crispr-2018}; closest-equivalent primary citation \texttt{paper:v6c03-cyranoski-ledford-2018}); the 2011--2012 Fouchier-Kawaoka H5N1 ferret-transmission gain-of-function controversy.}

\noindent
Synthetic biology applies engineering discipline to biological
systems: standard parts, characterised behaviour, hierarchical
abstraction, design-build-test-learn cycles. The reader who finishes
this chapter can specify a genetic circuit on paper, estimate its
robustness against parameter variation, name the principal failure
modes of CRISPR editing, and articulate the engineering and
ethical reasons that a containment plan precedes the experimental
plan. The field's frontier moves quickly. The chapter teaches the
durable engineering frame and uses 2010s and 2020s artefacts as
illustrative material rather than as durable specifications.

\section{What synthetic biology is}\pathtag{core}

Synthetic biology is the design and construction of biological
systems by engineering methods. The methods include: the use of
well-characterised standard parts; the design of new genetic
sequences and their introduction into cells; the prediction of
system behaviour from part-level specifications; the iterative
refinement of designs by measurement against the prediction. The
discipline emerged from molecular biology and genetic engineering in
the early 2000s, with three landmark papers in 2000 establishing
that genetic circuits could be designed and built that produced
prescribed dynamic behaviour: the Elowitz-Leibler oscillator
\cite[pp.~335--338]{paper:v6c09-elowitz-leibler-2000}, the Gardner-Cantor-Collins toggle switch \cite[pp.~339--342]{paper:v6c09-gardner-toggle-2000},
and the related early circuit demonstrations.

\engterm{Synthetic biology versus genetic engineering}. Genetic
engineering, dating from the 1970s, transferred genes between
organisms and modified specific genes for specific outcomes
(insulin-producing \emph{E.\ coli} in 1978, Roundup-Ready soybean
in 1996). Synthetic biology adds the engineering disciplines:
modularity (parts are reusable across designs), characterisation
(parts come with data sheets), abstraction (designs reason at the
device or system level, not at the nucleotide level), standardisation
(parts conform to interface specifications), and the explicit
design-build-test-learn workflow.

\engterm{The abstraction hierarchy}. Endy's foundational paper
\cite[pp.~451--452]{paper:v6c09-endy-foundations-2005} proposed a four-level
hierarchy: DNA (the substrate), parts (promoters, ribosome binding
sites, coding sequences, terminators), devices (parts wired together
to perform a function, like a logic gate or sensor), and systems
(devices composed into larger structures). The hierarchy works to
the extent that the abstractions hold: a part's behaviour depends
on its context (the cellular host, the metabolic state, the
neighbouring sequences), so the abstraction leaks. Managing the
leak is the field's central engineering problem.

\engterm{The four engineering challenges}. The field has not
achieved the reliable construction that Endy hoped for. The
challenges that remain
\cite[ch.~1]{text:v6c09-alon-systems-biology-2019}: context-dependent part
behaviour (a promoter that is strong in one strain is weak in
another); cellular resource competition (synthetic circuits draw
ribosomes, RNA polymerases, and amino acids away from the host's
own functions); evolutionary instability (synthetic constructs that
slow growth are selected against and lost on time scales of days);
and the absence of a quantitative theory that predicts circuit
behaviour from part specifications. The challenges have moderated
expectations: in 2026, a working synthetic biologist can build
toggle switches and oscillators in standard chassis on the time scale
of weeks, but the design of a non-trivial metabolic pathway still
requires substantial laboratory iteration.

\engterm{The field's two halves}. The field divides operationally
into circuit synthetic biology (toggle switches, oscillators, logic
gates, sensors) and metabolic engineering (redirecting a cell's
chemistry to produce a target molecule at industrial scale). The
circuits half cultivates the engineering ideals of modularity and
characterisation; the metabolic half is closer to chemical
engineering with a living catalyst. The two halves overlap in
fermentation, in bioreactor design, in the standard chassis
(\emph{E.\ coli}, \emph{S.\ cerevisiae}, CHO cells), and in the
practical pipeline from screen-to-pilot-to-production.

\section{Standard parts: BioBricks and the iGEM tradition}\pathtag{standard}

\input{volumes/06-life/ch09-synthetic-biology-metabolic-engineering/figures/fig-biobrick-plasmid}


The BioBricks project, initiated at MIT in 2003 by Tom Knight, Drew
Endy, and Randy Rettberg, formalised the part-level abstraction.
A BioBrick is a DNA sequence flanked by standard restriction sites
(\emph{EcoRI}-\emph{XbaI} on one side, \emph{SpeI}-\emph{PstI} on
the other) so that any two parts can be combined into a composite
part with the same flanking sites, allowing hierarchical assembly
\cite[\S\,2]{paper:v6c09-knight-biobricks-2003}. The standard enables
combinatorial design: a registry of 10{,}000 parts supports the
construction of many more composite parts without bespoke design at
each step.

\engterm{The Registry of Standard Biological Parts}. The Registry
hosts more than 20{,}000 BioBrick parts, contributed by participants
in the International Genetically Engineered Machine (iGEM)
competition since 2004. The Registry is a public catalogue: each part
has a unique identifier (BBa\_K725000, etc.), a sequence record, and,
in well-characterised cases, a data sheet with measured properties.
The dataset is heterogeneous. Parts contributed early are documented
better than parts contributed in the rush at competition deadlines;
parts in widespread use are characterised across contexts, while
single-use parts are not.

\engterm{The iGEM competition}. The iGEM competition (annual since
2004) is the world's largest synthetic-biology design contest. Teams
of undergraduates and high-schoolers design and build projects over
the summer and present them at the iGEM Jamboree. Past projects
include arsenic biosensors, dye-degrading bacteria, antibody-display
yeast, and a biological computer that performs counting. The
competition is the field's principal pipeline: a substantial
fraction of working synthetic biologists in industry and academia
came through iGEM. The competition's pedagogy is engineering
discipline applied to biology: planning, characterisation,
documentation, safety review. Whether the resulting parts and
devices are routinely reusable across other groups' projects is a
weaker claim; the Registry's reuse statistics show heavy use of a
small core (chassis, common reporters, common promoters) and very
little reuse of the long tail.

\engterm{Beyond BioBricks: Gibson assembly and Golden Gate}. The
flanking-restriction-site protocol of BioBricks is one assembly
method among several. Gibson assembly
\cite[pp.~343--345]{paper:v6c09-gibson-assembly-2009}, introduced in 2009,
uses a 5'-exonuclease, a polymerase, and a ligase in a single
isothermal reaction to fuse DNA fragments with overlapping ends. The
method removes the requirement for compatible restriction sites and
the scar sequences they leave. Golden Gate assembly, based on
type-IIS restriction enzymes that cut outside their recognition
sequence, enables seamless ordered assembly of multiple parts in a
single reaction. As of 2026, Gibson and Golden Gate have largely
displaced BioBricks for routine work; the BioBrick standard's
contribution is the conceptual one of a registry-based engineering
practice, rather than the specific assembly chemistry.

\section{Genetic circuits: switches, oscillators, logic}\pathtag{core}

A genetic circuit is a network of genes whose interactions produce
a prescribed temporal or logical behaviour. Three canonical circuits
illustrate the engineering discipline.

\input{volumes/06-life/ch09-synthetic-biology-metabolic-engineering/figures/fig-toggle-switch-state}

\engterm{The Hill function as the circuit engineer's workhorse}.
Every regulated promoter in the chapter's models obeys, to first
approximation, a Hill function. A repressor R that binds
cooperatively to its operator with apparent dissociation constant
$K$ and Hill coefficient $n$ shuts down transcription at a rate
\[
f(R) = \frac{1}{1 + (R/K)^{n}},
\]
so the protein production rate from a repressed promoter is
$\alpha_{\max} \cdot f(R) + \alpha_{\text{leak}}$, where
$\alpha_{\text{leak}}$ is the unrepressed background. For an
activator A acting cooperatively the function flips to
$g(A) = (A/K)^{n}/(1 + (A/K)^{n})$. The Hill coefficient $n$
captures the cooperativity of binding: $n = 1$ is a single binding
site (a graded response that responds to a 9-fold input change
between 10\% and 90\% output), $n = 2$ to $n = 4$ matches the
dimerised or tetramerised repressors actually used in synthetic
circuits (a sharper switch that responds within a 3-fold input
range), and $n \to \infty$ is the step response. Bistable circuits
require $n > 1$ for both repressors; oscillators require sufficient
cooperativity and time delay. The Hill function is a steady-state
approximation to the underlying receptor-binding equilibrium; the
working engineer treats it as the device's published transfer
function, and its parameters as data-sheet entries to be measured
and tabulated.

\engterm{Toggle-switch dynamics in detail}. The Gardner-Cantor-Collins
toggle has been analysed in detail because its two stable states are
the simplest case of designed bistability. The nullclines
$du/d\tau = 0$ and $dv/d\tau = 0$ are the sigmoidal curves
$u = \alpha_{1}/(1 + v^{n})$ and $v = \alpha_{2}/(1 + u^{m})$,
respectively. Their intersection determines the equilibria. For
$\alpha_{1} = \alpha_{2} = \alpha$ and $m = n$, symmetric analysis
shows the system has either one equilibrium (monostable, when
$\alpha \le \alpha_{c}$ with $\alpha_{c} = n/(n-1) \cdot (n-1)^{1/n}$
for the symmetric case; $\alpha_{c} \approx 2$ for $n = 2$ and
$\alpha_{c} \approx 1.75$ for $n = 3$), or three (bistable, with
two stable nodes on the diagonal $u = v$ symmetric counterparts and
a saddle at the symmetric point $u = v$). Bistability requires
$\alpha > \alpha_{c}$ and $n > 1$ simultaneously; either
non-cooperativity ($n = 1$) or weak production ($\alpha < \alpha_{c}$)
collapses the system to a single steady state. The original
Gardner construct used the lambda CI repressor and LacI (both
dimeric, $n \approx 2$) with strong promoters delivering
$\alpha \approx 5$ to $\alpha \approx 50$, so it operated well inside
the bistable regime
\cite[pp.~340--341]{paper:v6c09-gardner-toggle-2000}. The switching
inducer (heat for CI, IPTG for LacI) raised the effective dissociation
constant of one repressor's operator by a factor of 10 to 100, which
moved the corresponding nullcline outward in parameter space until
the relevant stable state lost its basin of attraction and the
system settled into the alternative state. The engineering virtue
of the toggle is that the inducer needs only act long enough to push
the system past the saddle; once across, the system completes the
transition on its own and remains in the new state after the
inducer is washed out.

\engterm{Repressilator dynamics in detail}. The Elowitz-Leibler
construct couples three repressor genes in a directed ring: protein
A represses gene $B$, B represses gene $C$, C represses gene $A$.
With dimensionless production rates $\alpha$ and a Hill coefficient
$n$ for each repressor, the symmetric model is
\begin{align*}
\frac{dp_{i}}{d\tau} &= -p_{i} + \frac{\alpha}{1 + p_{j}^{n}} + \alpha_{0},
\qquad (i, j) \in \{(1,3),\,(2,1),\,(3,2)\},
\end{align*}
where $\alpha_{0}$ is the leakage and $i$ refers to the protein
expressed from gene $i$. The symmetric equilibrium
$p_{1} = p_{2} = p_{3} = p^{*}$ exists for all parameters; linear
stability analysis around it shows the three eigenvalues are the
cube roots of $-X$ where
$X = n \alpha p^{*\,(n-1)} / (1 + p^{*\,n})^{2}$, scaled by the protein
decay rate. The symmetric equilibrium loses stability (Hopf
bifurcation) when one of the eigenvalues crosses the imaginary axis,
which occurs when $X > 8$ for $n = 2$ and $X > 2.6$ for $n = 4$
\cite[pp.~336--337]{paper:v6c09-elowitz-leibler-2000}. Strong
production rate $\alpha$ and sharp cooperativity $n$ favour
oscillation. The period of the limit cycle in the deterministic
limit is approximately $4\tau_{\text{prot}} \cdot \ln(\alpha)$ for
strong repression, where $\tau_{\text{prot}}$ is the protein half-life
including dilution. Elowitz and Leibler's 2000 strain produced an
oscillation period of roughly 150 minutes against a cell-division
time of 40 minutes; the mismatch made the oscillation noisy because
each daughter cell inherited a random share of the protein
inventory at division
\cite[p.~337]{paper:v6c09-elowitz-leibler-2000}. The 2016 refined
repressilator by Potvin-Trottier et al. engineered protein-decay
chains that brought the timescale closer to the cell cycle and
reduced cell-to-cell phase scatter
\cite[pp.~94--95]{paper:v6c09-potvin-trottier-repressilator-2016}.

\input{volumes/06-life/ch09-synthetic-biology-metabolic-engineering/figures/fig-repressilator}

\engterm{The toggle switch}. Gardner, Cantor, and Collins built the
canonical toggle in \emph{E.\ coli} in 2000
\cite[pp.~339--342]{paper:v6c09-gardner-toggle-2000}. Two repressor proteins
mutually inhibit each other's expression: protein A represses gene
$B$ (which expresses protein B); protein B represses gene $A$ (which
expresses protein A). The system has two stable states: A-high
B-low, and A-low B-high. A transient inducer (heat, an
exogenously added small molecule) can switch the system from one
state to the other; the system remains in the new state after the
inducer is removed. The model in dimensionless form is
\begin{align*}
\frac{du}{d\tau} &= \frac{\alpha_1}{1 + v^n} - u, \\
\frac{dv}{d\tau} &= \frac{\alpha_2}{1 + u^m} - v,
\end{align*}
where $u, v$ are the dimensionless protein concentrations, $\alpha_1,
\alpha_2$ are the dimensionless production rates, and $m, n$ are the
Hill cooperativities. The bistable regime requires $\alpha_1, \alpha_2
> 1$ and sufficient cooperativity ($m, n > 1$); for $m = n = 2$ and
balanced $\alpha$, the regime is centred on $\alpha \approx 3$. The
toggle is the synthetic-biology analogue of a flip-flop, and it
demonstrates that simple engineered networks can produce hysteresis.

\engterm{The repressilator}. Elowitz and Leibler built the canonical
oscillator in \emph{E.\ coli} in 2000
\cite[pp.~335--338]{paper:v6c09-elowitz-leibler-2000}. Three repressor
proteins arranged in a ring: A represses B, B represses C, C
represses A. The system has no stable equilibrium when the
repression is sufficiently strong; it oscillates with a period
determined by the protein production, decay, and dilution rates.
The original repressilator's period was about 150 minutes,
substantially longer than the bacterial cell cycle of 40 minutes,
so the oscillation was sloppy: phase drifted between cells, and
synchrony was rapidly lost. A refined repressilator
\cite[pp.~93--96]{paper:v6c09-potvin-trottier-repressilator-2016} of 2016
achieved better phase stability with engineered protein decay
chains. The repressilator is the synthetic-biology analogue of a
ring oscillator and demonstrates that simple engineered networks
can produce sustained limit cycles.

\engterm{Logic gates}. Boolean logic in DNA: AND gates produced
when two transcription factors must both bind to activate
transcription; OR gates when either of two activators suffices; NOT
gates by inversion through a repressor cascade. The gates compose
to produce non-trivial logic functions, though composition is
limited by the cross-talk that arises when too many gates share the
host's transcriptional machinery. In 2013, Stanford's Drew Endy and
collaborators reported a 16-input combinatorial gate in \emph{E.\
coli}; in 2017, MIT's Christopher Voigt and collaborators
demonstrated automated design and DNA synthesis of multi-gate
circuits (Cello, the genetic-circuit compiler). The engineering
direction is toward larger circuits with predictable behaviour; the
empirical record shows that prediction degrades rapidly with
circuit size.

\engterm{Adaptive and feedback circuits}. Negative feedback in a
genetic circuit produces homeostasis (the output is held constant
against perturbations), adaptation (the output responds to a step
input and then returns to baseline), or oscillation (when the
feedback is delayed). Alon's textbook
\cite[chs.~3--5]{text:v6c09-alon-systems-biology-2019} develops the canonical
design principles: integral feedback for robust homeostasis,
incoherent feed-forward loops for fold-change detection, coherent
feed-forward loops for noise filtering. The principles transfer
from naturally occurring regulatory networks (which evolution has
already optimised for these problems) to designed networks (which
the engineer constructs from the same building blocks).

\begin{archetype}[Control: feedback in genetic circuits]
A genetic circuit is a small dynamical system with the same control
structure as any other: state variables (protein concentrations),
dynamics (production minus degradation, with regulatory terms
depending on the state), inputs (inducers, environmental
conditions), and outputs (reporter concentrations). The engineer
designs the regulatory architecture (feedback, feed-forward, cascade,
oscillator) to achieve the target behaviour, and characterises the
result by measurement. The discipline is control engineering with a
particularly slow plant and particularly noisy actuators.
\end{archetype}

\begin{estimation}
\textbf{Question.} Estimate the steady-state protein concentration
in a bacterial cell expressing a synthetic gene from a moderate
promoter. Assume the gene is transcribed at 10 mRNAs per cell per
generation, each mRNA is translated 30 times before being degraded,
the protein has degradation rate set by cell division (one division
per 40 minutes), the cell volume is $1\,\si{\micro\meter}^3$.

\smallskip
\textit{Estimate before reading on.}

\smallskip
The protein production rate is 10 mRNA $\times$ 30 protein per mRNA
$= 300$ proteins per cell per generation, or
$300/(40\,\text{min}) = 7.5$ proteins per minute per cell. The
protein degradation rate, if dilution-only, is $\ln 2/40\,\text{min}
= 0.017\,\text{min}^{-1}$. At steady state, production equals
degradation times concentration: $7.5 = 0.017\,N$, so $N = 440$
proteins per cell. The concentration is $440 / (1\,\si{\micro\meter}^3
\cdot 10^{-15}\,\text{L}/\si{\micro\meter}^3) = 440/(6 \times 10^{8})
\approx 7 \times 10^{-7}\,\text{M} = 0.7\,\si{\micro\text{M}}$.

\smallskip
\textit{Postmortem.} The order of magnitude (a few hundred to a few
thousand proteins per cell, $\sim 1\,\si{\micro\text{M}}$) matches
the measured abundance of moderately expressed proteins in
\emph{E.\ coli} \cite[ch.~2]{text:phillips-physical-biology}. A
strong inducible promoter (T7, $P_{\text{Lac}}$ fully induced) raises
production by 10--100 times and saturates the cell's translation
machinery, producing the typical 10{,}000--100{,}000 proteins per
cell. A weak promoter (PoPS in the BioBrick Registry's low range)
lowers production by 10--100 times and produces dozens of proteins
per cell. The two-orders-of-magnitude range is the design range a
genetic-circuit engineer works within; the estimate above is the
midpoint that an initial design defaults to.
\end{estimation}

\begin{estimation}
\textbf{Question.} Estimate the off-target probability for a single
20-nucleotide Cas9 guide against the human genome. The genome is
$3.2 \times 10^{9}$ base pairs, the PAM constraint reduces accessible
sites by a factor of about 16 (the NGG PAM means about 1 in 8 of
both strands, $1/16$ of all sites genome-wide), and Cas9 tolerates
single mismatches at most positions of the guide with about 10\%
of on-target efficiency, two mismatches at about 1\% in
favourable PAM-distal positions, and three mismatches only at
0.01\% in select positions.

\smallskip
\textit{Estimate before reading on.}

\smallskip
With no mismatch tolerance, the expected number of perfect-match
sites in the genome is $(3.2 \times 10^{9}/16)/4^{20} = 2 \times 10^{8}/10^{12} \approx 2 \times 10^{-4}$;
a randomly chosen 20-mer with a PAM matches no other genomic site
with high probability. The off-target risk arises from
mismatch-tolerant binding. Sites with one mismatch in PAM-distal
positions: there are $20 \cdot 3$ single-mismatch variants of the
guide, each appearing in the genome roughly
$(3.2 \times 10^{9}/16)/4^{19} = 6 \times 10^{-4}$ times in expectation,
times 10\% binding efficiency, gives an expected
$20 \cdot 3 \cdot 6 \times 10^{-4} \cdot 0.1 \approx 4 \times 10^{-3}$
single-mismatch off-target events per edit per cell. Sites with two
mismatches: $\binom{20}{2} \cdot 3^{2} \cdot 4 \times 10^{-3} \cdot 0.01 \approx 0.07$ per cell. The two-mismatch term dominates.
For a population of $10^{8}$ edited cells, the expected
two-mismatch off-target hit count is $\sim 7 \times 10^{6}$. The
mitigations (genome-wide guide selection, high-fidelity Cas9
variants, paired nickases, base editors) reduce these rates by
two to four orders of magnitude before clinical use.

\smallskip
\textit{Postmortem.} Hsu et al. (2013) and Pattanayak et al. (2013)
report the empirical off-target distributions from genome-wide
binding assays
\cite[pp.~828--829]{paper:v6c09-hsu-offtarget-2013}\cite[pp.~840--843]{paper:v6c09-pattanayak-offtarget-2013}. The
order-of-magnitude estimate above accounts for the broad
distribution: most edits produce zero detected off-targets when
guides are selected by CRISPRoff or CHOPCHOP scoring, but
unscreened guides routinely yield several off-targets at low
frequency. The clinical applications (Casgevy, exa-cel) selected
guides that score well by the genome-wide off-target search and
confirmed off-target absence by GUIDE-seq, CIRCLE-seq, or whole-genome
sequencing of the edited cell product before regulatory submission.
The estimation establishes the order of magnitude that clinical
practice must beat by design and characterisation, not by hope.
\end{estimation}

\section{Metabolic engineering: redirecting cellular flux}\pathtag{core}

A cell is a chemical factory. Its substrate is glucose (or another
carbon source); its products are biomass, CO$_2$, and a few
metabolic by-products. Metabolic engineering redirects the cell's
chemistry to produce a target molecule at engineered yield, at
engineered rate, in an engineered host.

\engterm{Flux balance analysis}. The cell's metabolic network is a
graph: nodes are metabolites, edges are enzyme-catalysed reactions.
At steady state, each metabolite's concentration is constant, so
the flux into it equals the flux out:
\[
\mathbf{S} \cdot \mathbf{v} = \mathbf{0},
\]
where $\mathbf{S}$ is the $m \times n$ stoichiometric matrix
($m$ metabolites, $n$ reactions) and $\mathbf{v}$ is the flux
vector. The system is underdetermined (more reactions than
metabolite balances). The remaining degrees of freedom are closed
by an objective function (typically biomass production) and by
bounds on individual fluxes (substrate uptake, maximum enzyme
capacities, irreversibility constraints). The resulting linear
program is flux balance analysis (FBA)
\cite[chs.~3--4]{text:v6c09-palsson-systems-biology-2015}; commercial FBA
software (COBRA Toolbox, OptKnock, OptStrain) solves the LP and
identifies the gene knockouts and over-expressions that redirect
flux toward a desired product. FBA predicts cellular growth rates
within a factor of two for \emph{E.\ coli} on minimal media; it
predicts the qualitative effects of gene knockouts correctly more
often than not; it gets wrong quantitative predictions when the
enzyme kinetics, the regulation, or the cellular crowding deviates
from the LP's assumptions.

\engterm{Three classical metabolic-engineering successes}. The field
has produced reliable industrial products. Three named cases:

\begin{itemize}[leftmargin=*]
\item \emph{Recombinant insulin}. Genentech and Eli Lilly engineered
  \emph{E.\ coli} to produce human insulin in 1978; Humulin reached
  market in 1982 \cite[ch.~1]{text:v6c09-demain-industrial-microbiology-2017}.
  The product replaced bovine and porcine insulin for diabetes
  treatment. The engineering challenge was to produce a small,
  cysteine-rich, two-chain protein in a bacterial host that cannot
  perform the post-translational disulfide formation; the solution
  was to express the two chains separately, denature, and refold them
  \emph{in vitro}.
\item \emph{Recombinant artemisinin}. Jay Keasling's group at UC
  Berkeley engineered yeast to produce the precursor amorphadiene,
  followed by chemical conversion to artemisinin (the antimalarial)
  \cite[pp.~528--532]{paper:v6c09-paddon-artemisinin-2013}. The route
  expressed plant-derived enzymes in yeast and tuned the upstream
  isoprenoid pathway; Sanofi commercialised it in 2013--2014. The
  programme illustrated the time scale: roughly a decade from
  proof-of-concept to commercial production, with multiple iterations
  on host, pathway, and process.
\item \emph{Recombinant 1,3-propanediol}. DuPont and Genencor (now
  IFF) engineered \emph{E.\ coli} to produce 1,3-propanediol (a
  monomer for the Sorona polymer) from glucose
  \cite[ch.~9]{text:v6c09-demain-industrial-microbiology-2017}. The
  programme combined glycerol-producing yeast pathways and
  glycerol-to-propanediol bacterial pathways into a single
  \emph{E.\ coli} host. Commercial production started in 2006.
\end{itemize}

The three programmes share an engineering pattern: a high-value
target (drug or specialty chemical); a host chosen for its known
fermentation properties; pathway redesign by FBA-guided
identification of bottlenecks and rate-limiting enzymes; and
several iterations of mutagenesis, directed evolution, and process
optimisation before a commercially viable strain is reached.

\input{volumes/06-life/ch09-synthetic-biology-metabolic-engineering/figures/fig-metabolic-flux}

\engterm{Yield, titer, and productivity}. The three numbers a
metabolic-engineering programme reports. \engterm{Yield} is grams of
product per gram of substrate; \engterm{titer} is grams of product
per litre of broth; \engterm{productivity} is grams of product per
litre per hour. Industrial targets typically demand yield greater
than 0.4 g/g (for a non-trivial target like artemisinin), titer
greater than 50 g/L, and productivity greater than 1 g/L/h. A
laboratory-scale demonstration that hits one of the three may still
be a decade from a commercially viable process. The chapter on
bioreactors (Chapter~\ref{vol06:ch10}) develops the engineering of
the next-stage scale-up.

\section{CRISPR and gene editing engineering-first}\pathtag{core}

The CRISPR-Cas9 system, derived from a bacterial adaptive-immunity
mechanism, enables targeted modification of DNA in living cells. The
basic system: a single guide RNA (sgRNA) of about 20 nucleotides
matching the target DNA sequence; the Cas9 nuclease, which uses the
sgRNA to find and cut both strands of the target DNA; the cell's
own DNA-repair machinery, which closes the break either by
non-homologous end joining (NHEJ, error-prone) or by homology-directed
repair (HDR, accurate when a template is provided). Doudna and
Charpentier's 2012 paper \cite[pp.~816--821]{paper:v6c09-jinek-crispr-2012}
demonstrated programmable cleavage in vitro; Zhang and Church
demonstrated the same in mammalian cells in 2013. The 2020 Nobel
Prize in Chemistry was awarded to Doudna and Charpentier.

\engterm{The basic editing operation}. (1) Design an sgRNA matching a
20-nucleotide target adjacent to a protospacer-adjacent motif (PAM,
typically NGG for the \emph{S.\ pyogenes} Cas9). (2) Deliver Cas9
protein and sgRNA into the cell, either as DNA, mRNA, or
ribonucleoprotein. (3) Optionally deliver a repair template if HDR
is desired. (4) Screen the resulting cells for the desired edit.
The screening is necessary: the editing efficiency is typically
10\%--90\% for NHEJ, 0.1\%--10\% for HDR, with wide variation by
target site, cell type, and delivery method.

\input{volumes/06-life/ch09-synthetic-biology-metabolic-engineering/figures/fig-crispr-mechanism}

\engterm{Off-target effects}. Cas9 cuts not only the intended target
but also sites with imperfect sequence match. Off-target cutting is
the field's principal engineering challenge for clinical
applications. The mitigations: high-fidelity Cas9 variants (eSpCas9,
SpCas9-HF1, HypaCas9), tighter PAM-recognition variants, paired
Cas9-nickases that require two adjacent cuts for a double-strand
break, base editors that change a single base without cutting, prime
editors that write short sequences without cutting. Each approach
trades efficiency for specificity; for a clinical application, the
specificity must be high enough that the off-target rate is below
the spontaneous mutation rate of the cells being edited (typically
$10^{-8}$ per site per generation), a stringent target that current
tools meet only at carefully selected target sites.

\engterm{Off-target rate by sgRNA design}. The off-target frequency
depends sharply on the position of the mismatch within the 20-nt
guide. Hsu et al. (2013)\cite[pp.~828--829]{paper:v6c09-hsu-offtarget-2013}
measured the genome-wide binding profile of Cas9 by ChIP-seq across
hundreds of guide-target pairs and reported three rules. First, the
PAM-proximal 8--12 nt (the ``seed'' region) is intolerant of
mismatches; a single mismatch in this region typically drops binding
by at least an order of magnitude. Second, PAM-distal mismatches
(positions 1--8 counted from the PAM-distal end) are tolerated:
single mismatches preserve about 10\% of on-target activity, and
double mismatches separated by at least 4 positions preserve about
1\%. Third, transition mismatches (purine-to-purine or
pyrimidine-to-pyrimidine substitutions) are more readily tolerated
than transversions. Pattanayak et al. (2013)
\cite[pp.~840--843]{paper:v6c09-pattanayak-offtarget-2013}
independently profiled off-target activity by an in vitro library
selection and reproduced the same seed-distal asymmetry, with an
overall off-target frequency at sites with three or more mismatches
in the seed region below their detection floor of $10^{-5}$. The
design implications: pick a target sequence with no genome-wide
sites containing at most three PAM-distal mismatches; verify by
genome-wide assay (GUIDE-seq, CIRCLE-seq) at the cell type of
interest; if the patient population includes variants near the
target site, characterise the off-target profile in each variant
background separately.

\engterm{Other failure modes uncovered by clinical translation}.
The 2018--2019 wave of CRISPR-Cas9 clinical preparation surfaced
failure modes that off-target screening had not anticipated. Ihry
et al. (2018)\cite[pp.~939--942]{paper:v6c09-ihry-p53-2018} showed
that Cas9-induced double-strand breaks activate the p53 DNA-damage
response in human pluripotent stem cells, selecting for p53-deficient
cells in any edited population; an edited cell product enriched in
p53-deficient cells carries a long-term cancer risk that no
off-target screen detects. Cellini et al. (2018)
\cite[pp.~767--770]{paper:v6c09-kosicki-deletions-2018} reported
that Cas9 editing produces large deletions (kilobase scale, sometimes
megabase scale) at the target site at a higher frequency than the
small indels for which most assays test, and that translocations
between the target and off-target sites occur at measurable rates
when both are cut in the same cell. The clinical lesson: an edit
product validated by amplicon-sequencing of the target site alone
is not validated; long-range sequencing, karyotyping, and p53-status
assays became standard before regulatory submission of any
Cas9-edited cell product after 2019.

\engterm{CRISPR's first clinical use}. The first regulatory approvals
came in 2023. Casgevy (exa-cel), developed by Vertex and CRISPR
Therapeutics, treats sickle-cell disease and beta-thalassemia by
disrupting the BCL11A gene in the patient's hematopoietic stem cells
ex vivo, then reinfusing the cells. The disrupted gene allows fetal
hemoglobin to be expressed, partially compensating for the defective
adult hemoglobin. The therapy was approved by the FDA in December
2023 \cite{web:v6c09-fda-casgevy-2023}. The engineering frame: a
single-organ target (the bone marrow); an autologous, ex vivo
procedure that limits in vivo delivery problems; a payoff that
justifies the high cost ($2.2$ million USD per patient, as of 2024
pricing). Subsequent approved CRISPR therapies will have to meet
similar thresholds on target specificity, delivery, and clinical
benefit.

\engterm{The He Jiankui case}. In November 2018, the Chinese
biophysicist He Jiankui announced that he had used CRISPR to edit
the CCR5 gene in two human embryos, resulting in the birth of
twin girls. The announcement triggered immediate condemnation from
the international scientific community. The work violated
multiple norms: clinical research on heritable human editing was
prohibited in China and elsewhere; the targeted gene was not a
disease-causing variant (CCR5 is a co-receptor for HIV entry, and
its disruption protects against HIV, but the parents were not at
high risk); the consent process was inadequate; the off-target
rate was not characterised. He was tried in China in 2019 and
sentenced to three years in prison. The case is the field's
canonical breach: clinical use of heritable editing without the
engineering, regulatory, or ethical scaffolding that any such use
requires \cite[pp.~127--132]{text:v6c09-nas-heritable-editing-2020}.

\section{Directed evolution as a design tool}\pathtag{standard}

When the design problem is too complex for rational engineering,
directed evolution substitutes selection for design. Frances
Arnold's group at Caltech demonstrated the method on enzyme
engineering through the 1990s and 2000s; Arnold shared the 2018
Nobel Prize in Chemistry for the programme
\cite[pp.~1747--1750]{paper:v6c09-arnold-directed-evolution-1998}. The
method: take a starting protein with weak activity for a desired
function; subject its gene to mutagenesis (error-prone PCR, DNA
shuffling, structure-guided mutation); express the mutant library
in a host; screen or select for activity; retain the best
performers; iterate.

\engterm{The selection pressure must do the work}. The library
contains many variants ($10^4$ to $10^{10}$ per round); the
selection identifies the few with improved activity. A good
selection is the difference between success and failure. A bad
selection (one that does not actually penalise lack of the target
function, or that selects for an off-target trait) leads to
``cheating'' solutions that do not generalise. A high-throughput
screen (FACS for cells; microtitre-plate enzyme assays) supports a
larger library than a low-throughput one (NMR of purified protein,
one variant at a time).

\engterm{Engineering for evolution}. The library design matters as
much as the selection. Saturation mutagenesis at a small number of
well-chosen positions (informed by structure and homology) reaches
useful improvements faster than random mutagenesis across the whole
gene. DNA shuffling \cite[pp.~389--391]{paper:v6c09-stemmer-shuffling-1994}
recombines mutations from many parental genes, allowing the
selection to find combinations of mutations that none of the
parents possessed.

\engterm{Phage-assisted continuous evolution (PACE)}. The discrete
round of directed evolution (library construction, selection,
sequencing, library reconstruction) takes one to two weeks of
bench work per generation. PACE, introduced by Esvelt, Carlson, and
Liu in 2011 \cite[pp.~499--503]{paper:v6c09-esvelt-pace-2011},
collapses the round into a continuous biological process. The gene
of interest is encoded on an M13 bacteriophage with the gene III
coat protein replaced by the evolution target; the phage circulates
through a chemostat (the ``lagoon'') containing \emph{E.\ coli} that
express a fitness-coupled gene III protein only when the target
gene's product carries the desired activity. Phages whose target
genes produce active product produce more progeny; phages whose
target genes are inactive produce dilute, infectious-incompetent
progeny that wash out of the lagoon at the chemostat dilution rate.
Mutagenesis is supplied by a small-molecule-inducible mutator
plasmid (typically the MP6 plasmid expressing error-prone DNA
polymerase variants under arabinose control). A single PACE run
performs the equivalent of several hundred discrete rounds of
directed evolution per week in a tabletop setup, with the
selection pressure tunable by the chemostat dilution rate and the
mutation rate tunable by the arabinose concentration. PACE has
been used to evolve T7 RNA polymerase variants with altered
promoter recognition, base editors with shifted target preferences,
and Cas variants with relaxed PAM requirements, in projects whose
non-PACE-equivalent timeline would have been a graduate career.

\input{volumes/06-life/ch09-synthetic-biology-metabolic-engineering/figures/fig-pace-pipeline}

\engterm{Modern directed evolution: machine learning in the loop}.
Since 2020, directed evolution has been augmented by ML models
trained on the variant-activity data already accumulated. A model
trained on the first few rounds of a directed-evolution programme
can predict which untested variants are likely to be improved,
guiding the design of the next round and reducing the experimental
burden by 10--100$\times$ in favourable cases. The combination of
directed evolution and ML is the standard contemporary practice in
enzyme engineering; the chapter on bioinformatics
(Chapter~\ref{vol06:ch11}) develops the ML methods.

\section{Containment and biosafety levels}\pathtag{core}

The work of synthetic biology is performed in physical and
biological containment. The two are layered: physical containment
(laminar-flow hoods, sealed centrifuges, autoclaves, locked
freezers) restricts the release of organisms; biological containment
(host strains that cannot survive outside the laboratory, dependence
on engineered nutrients, kill switches that activate on escape)
restricts the survival of any organism that does escape. The
discipline is the same that microbiology and infectious-disease
research has developed over a century; synthetic biology inherits
it and extends it for engineered organisms whose hazard profile may
differ from natural strains.

\input{volumes/06-life/ch09-synthetic-biology-metabolic-engineering/figures/fig-bsl-containment}

\engterm{The four biosafety levels (BSL)}. The classification scheme
from the CDC and WHO categorises laboratories by the agents they
handle \cite[\S\,II]{web:v6c09-cdc-bmbl-6-2020}: BSL-1 for non-pathogenic
organisms (most laboratory \emph{E.\ coli} K-12 derivatives); BSL-2
for organisms that cause moderate disease in healthy adults via
ingestion or percutaneous exposure (\emph{Staphylococcus aureus},
\emph{Salmonella}); BSL-3 for organisms transmitted by aerosol that
cause serious disease (\emph{Mycobacterium tuberculosis}, SARS-CoV-2);
BSL-4 for organisms with high mortality and no available treatment
(Ebola, Marburg, smallpox). Each level specifies physical
containment, personal protective equipment, training requirements,
and the institutional oversight required to operate. The
classification of a synthetic-biology project is determined by the
hazard of the chassis, of the inserted genes, and of any reasonably
foreseeable misuse of the engineered organism.

\engterm{BSL-1 through BSL-4 in concrete terms}. The CDC-NIH
\emph{Biosafety in Microbiological and Biomedical Laboratories}
manual (6th edition, 2020)\cite[\S\,II]{web:v6c09-cdc-bmbl-6-2020}
codifies the requirements. BSL-1 specifies an ordinary laboratory:
work on open benches with standard personal protective equipment
(lab coat, gloves, eye protection); decontamination of surfaces and
waste with standard disinfectants; no special engineering controls.
Typical work: laboratory \emph{E.\ coli} K-12 derivatives, baker's
yeast, soil bacteria not known to cause disease in healthy adults.
BSL-2 adds restricted access to the laboratory during work, a
biological safety cabinet for procedures that generate aerosols,
sharps controls, and a written biosafety manual; the laboratory
must train workers on the specific agents being handled. Typical
work: HIV in cell culture, \emph{Salmonella} or
\emph{Staphylococcus aureus} cultures, most human pathogens that
require ingestion or direct contact for transmission. BSL-3 adds
controlled-access entry through an anteroom; a directional airflow
that pulls air from the corridor into the laboratory and from the
laboratory through HEPA filters before exhaust; all procedures
performed in a biological safety cabinet; respiratory protection
for personnel; sealed penetrations through walls and floors;
medical surveillance and serum banking of workers; and an
institutional biosafety committee that approves protocols before
work begins. Typical work: \emph{Mycobacterium tuberculosis},
SARS-CoV-2, yellow-fever virus, Rift Valley fever virus, prion
proteins from human or animal-prion sources. BSL-4 adds positive-pressure
``space suits'' supplying breathing air independent of the room
atmosphere, double-door autoclaves and dunk tanks for
decontamination of all material leaving the laboratory, shower-in
and shower-out exits, redundant power and HVAC for fail-safe
containment, and 24/7 security monitoring. Typical work: Ebola
virus, Marburg virus, smallpox virus (held only at the CDC and
Vector Institute under WHO authority), Lassa fever virus. The
2020 census of US BSL-4 facilities counted 15 facilities; the
global count was around 60. The escalation between levels is
non-linear in cost: a BSL-2 retrofit of an existing BSL-1
laboratory costs on the order of \$30{,}000 in equipment and
training; a BSL-3 laboratory requires dedicated space and HVAC
and runs to several million dollars in construction; a BSL-4
suite runs to \$50--200 million dollars in construction. The cost
gradient is itself an engineering control: it concentrates the
highest-hazard work in a small number of facilities, where the
discipline of operations can be maintained and audited.

\engterm{Built-in containment}. Some engineered organisms include
auxotrophic dependencies (the strain cannot grow without a synthetic
amino acid not present in nature), conditional lethality (a
kill switch triggered by a sensor for an environmental condition the
laboratory is engineered to avoid), or recoded genetic codes (the
organism reads its DNA differently from any natural organism, so
horizontal gene transfer to natural strains produces non-functional
products). Each provides a layer of biological containment.
Empirical studies of escape frequency report rates of $10^{-8}$ to
$10^{-12}$ per cell per generation for engineered containment
strategies, depending on the design and the selection pressure for
escape mutants \cite[pp.~115--118]{paper:v6c09-rovner-genetic-code-2015}.

\engterm{Dual-use research of concern (DURC)}. Some research has
clear scientific or public-health value and clear potential for
misuse. The most prominent recent case: gain-of-function research on
H5N1 influenza by Ron Fouchier's and Yoshihiro Kawaoka's groups in
2011--2012, which engineered mutants of H5N1 transmissible by
respiratory droplets between ferrets (a model organism for human
transmission) \cite[pp.~534--541]{paper:v6c09-imai-h5n1-2012}. The studies
were intended to identify mutations that could naturally accumulate
and produce a pandemic; the publication of the results raised the
risk of replication by malicious actors. The US government imposed a
moratorium on funding GoF research in 2014, lifted in 2017 with a
revised framework. The framework is still contested: the engineering
question is whether the prediction value of the research exceeds the
risk it generates; reasonable engineers disagree.

\section{Failure: gene-edit off-target, gain-of-function, dual-use}\pathtag{core}

Synthetic biology fails in characteristic ways. The chapter's
failure section names three.

\engterm{Gene-edit off-target effects}. CRISPR-Cas9 cuts at sites
whose sequence matches the sgRNA imperfectly. The off-target
frequency depends on the sequence similarity, on the chromatin
context, on the duration of Cas9 expression, and on the cell type.
For clinical applications, off-target cutting can produce unintended
mutations in cancer-driving genes (tumour suppressors, proto-oncogenes)
or in essential genes. The mitigations (high-fidelity variants, base
editors, prime editors, careful target selection) reduce but do not
eliminate the rate. The failure mode is engineering with an editing
tool whose precision is insufficient for the application; the
response is to characterise the tool's precision in the relevant
context and to validate that the off-target rate is acceptable for
the intended use.

The 2018 wave of high-profile clinical-preparation studies (Hsu
2013; Pattanayak 2013; Cellini/Kosicki 2018; Ihry 2018) shifted the
field's understanding of what counts as ``off-target.''
Cellini-Kosicki et al. (2018)
\cite[pp.~767--770]{paper:v6c09-kosicki-deletions-2018} reported
that target-site cleavage produces large structural variants
(deletions of kilobases to megabases, chromosomal translocations,
and inversions) at a frequency the small-indel amplicon assays had
not detected, raising the burden of validation for any
Cas9-edited cell product. Ihry et al. (2018)
\cite[pp.~939--942]{paper:v6c09-ihry-p53-2018} reported that
Cas9-induced double-strand breaks trigger the p53 damage response
and reduce the survival of cells carrying intact p53 by an
order of magnitude relative to p53-deficient cells in the same
population. An edit performed on a cell population with normal
p53 surveillance enriches the surviving cells for p53 loss-of-function,
a tumour-suppressor knockout the protocol never intended to make.
The He Jiankui case (registry entry
\texttt{acc:he-jiankui-crispr-2018}) closed the loop: the embryos
were mosaic with novel insertions and deletions at CCR5, the
off-target characterisation was inadequate, and the children's
long-term cancer risk was not assessed
\cite[pp.~127--132]{text:v6c09-nas-heritable-editing-2020}.

\engterm{Gain-of-function and dual-use risk}. Some experiments
generate organisms more dangerous than any natural strain. The
2011--2012 ferret-transmission H5N1 work by Fouchier (Erasmus) and
Kawaoka (Wisconsin/Tokyo) is the canonical case; the 2014 US
funding moratorium and the 2017 P3CO framework are the canonical
regulatory response\cite[pp.~534--541]{paper:v6c09-imai-h5n1-2012}.
Imai et al. (2012) engineered H5N1 mutants transmissible by
respiratory droplets between ferrets (the model organism for human
transmission), identifying four mutations (two in the haemagglutinin
receptor-binding site, one in HA stalk, one in PB2) that
collectively conferred airborne transmissibility in the model.
The pre-publication safety review took six months; the published
results raised the global biosecurity concern that the same
mutations could be introduced into circulating H5N1 strains by an
actor with modest molecular-biology capability. The US National
Science Advisory Board for Biosecurity (NSABB) initially recommended
that the manuscripts be redacted; the recommendation was withdrawn
after a WHO consultation, and the papers were published in 2012.
The October 2014 White House Office of Science and Technology
Policy moratorium paused federal funding for gain-of-function
research on influenza, SARS, and MERS pending an HHS-led risk
assessment. The December 2017 HHS P3CO (Potential Pandemic
Pathogen Care and Oversight) framework restored funding under a
case-by-case review that requires demonstrating the research's
scientific value relative to its biosafety and biosecurity risks.
The framework has been criticised both for over-restriction (a
2014 NIH-funded SARS chimaeric study had to be paused pending
P3CO review) and for under-restriction (the 2015 Baric-Menachery
SARS-like chimera study was approved under a transitional review
that critics argued did not meet P3CO standards).

The DURC (dual-use research of concern) framework, codified in
2012 US federal policy and revised in 2024 to cover a broader set
of pathogens and experimental categories, applies to research that
could be ``directly misused to pose a significant threat with broad
potential consequences to public health and safety''
\cite[\S\,III]{web:v6c09-hhs-durc-2024}. The seven experimental
categories DURC covers: enhancing harmful consequences of an agent;
disrupting effective immunity to an agent; conferring resistance to
useful therapeutic interventions; increasing transmissibility of an
agent; altering host range or tropism; enabling evasion of
diagnostic detection; enabling weaponisation. The framework requires
institutional review before submission of grant applications and
publication of results; it does not prohibit any specific
experiment, but it requires the institution and the investigator
to identify and mitigate the misuse risk. The 2024 revision
broadened the scope of covered pathogens (to include any pathogen
with pandemic potential) and tightened the institutional review
requirements; reasonable engineers and policymakers disagree on
whether the revision strikes the right balance between scientific
progress and biosecurity. The 2026 state of the framework is
incomplete.

\engterm{Evolutionary instability of engineered constructs}. A
synthetic circuit that imposes a fitness cost on its host is
selected against in laboratory culture. Constructs lose function on
time scales of days to weeks: a strong promoter accumulates
inactivating mutations, an antibiotic-resistance marker is lost when
the antibiotic is not present, a metabolic pathway is downregulated
by host regulation. The engineering responses: minimise the fitness
cost by tuning expression levels; use auxotrophic markers that
maintain selection (the strain cannot grow without the plasmid);
integrate constructs into the chromosome where they are more stable
than on plasmids; design redundant or self-stabilising circuits.
The failure mode is to deploy a circuit in a strain that has not
been characterised for stability over the time scale of the intended
use; the response is to characterise stability empirically as a
design specification.

\engterm{He Jiankui as ethics-and-safety failure}. The He Jiankui
case (November 2018; registry entry
\texttt{acc:he-jiankui-crispr-2018}) is the field's canonical
ethics-and-safety failure. The technical failure: the editing
in human embryos did not reproduce the natural CCR5-$\Delta$32
HIV-protective allele but produced novel deletions of 15 and 4
base pairs in one child and a 1-base insertion in the other; both
embryos were mosaic, with different cells carrying different edits,
indicating Cas9 activity continued past the first cleavage division;
off-target characterisation was not performed at clinical-grade
sensitivity\cite{paper:v6c03-musunuru-2019}. The organisational
failure: He bypassed Chinese national regulation explicitly
prohibiting heritable embryo editing; the institutional ethics
committee documentation was partly forged; informed consent did not
adequately describe the experimental nature of the work or its
risks. The clinical failure: the parents faced near-zero vertical
HIV transmission risk under standard sperm-washing protocols, so
the supposed benefit did not justify the intervention; the
children's long-term cancer risk from Cas9-induced large
deletions and from any p53 selection was not assessed; no plan
existed for long-term follow-up of the children. The regulatory
response: the Chinese government convicted He of illegal medical
practice in December 2019 and sentenced him to three years'
imprisonment; the WHO Expert Advisory Committee on Developing
Global Standards for Governance and Oversight of Human Genome
Editing issued binding global standards in 2021; the International
Commission on the Clinical Use of Human Germline Genome Editing
(2020) established the international consensus that heritable
human genome editing should not proceed until standards for safety,
efficacy, and societal authorisation are in place. The engineering
lesson: a single investigator at a single institution with access
to in vitro fertilisation infrastructure can, given sufficient
determination, conduct heritable genome editing without national
approval. Norms-without-enforcement do not constrain a determined
defector. The regulatory framework that prevents the next case
must therefore be binding and enforceable, not merely
recommendatory\cite[pp.~127--132]{text:v6c09-nas-heritable-editing-2020}.

\begin{principle}[Synthetic biology inherits its risks from natural biology and adds its own]
The risks of synthetic biology are the risks of natural biology
(pathogenicity, transmissibility, host range, environmental impact)
applied to organisms whose biology has been deliberately modified.
The modification can reduce risk (by removing pathogenicity factors,
by limiting environmental fitness) or increase it (by adding
factors, by extending host range). The engineering responsibility
is to characterise the modification's effect on the risk before the
modification is deployed, and to design layered containment that
covers the residual risk. The historical record of natural-history
risks (introduced species, antibiotic resistance, viral spillover)
provides the base rate; the engineering record of designed
containment (BSL classifications, kill switches, recoded codes)
provides the mitigation.
\end{principle}

\section{Project}\pathtag{core}

\begin{project}[Design a genetic toggle switch on paper]
The reader designs a genetic toggle switch in \emph{E.\ coli} and
defends its predicted robustness against parameter variation.

\textbf{Track}. Paper design and simulation.
\textbf{Hazard class}: low (no laboratory work).

\textbf{Specification}.

\begin{enumerate}[leftmargin=*]
\item Identify two repressor proteins and the promoters they
  regulate. Justify the choice from data sheets in the Registry of
  Standard Biological Parts (or from primary literature). Specify
  the ribosome binding site for each, the protein degradation tag
  if used, and the plasmid backbone or chromosomal integration site.
\item Write down the dimensional equations for the time evolution of
  the two repressor concentrations. State the production rate,
  degradation rate, Hill cooperativity, and any leakage. Identify
  the two stable equilibria.
\item Solve the dimensionless model numerically. Plot the phase
  portrait. Identify the basin of attraction of each stable state
  and the unstable saddle separating them.
\item Sweep each parameter (production rate, degradation rate, Hill
  cooperativity) by $\pm 50\%$ and report which parameters remove
  bistability. Document the robust region of parameter space.
\item Identify the inducer for each state-switching event.
  Specify the inducer concentration and duration required to switch
  state.
\item Identify three failure modes (one for each of: protein
  degradation tag failure, leakage at the promoter, growth-rate
  dependence of the model). For each, name a measurable signature
  and propose a corrective design.
\end{enumerate}

\textbf{Deliverable}. A 10- to 15-page report with a circuit diagram,
parameter table, simulation results, robustness analysis, and
failure-mode discussion.

\textbf{Time}. 10--20 hours of analysis and writing.
\end{project}

\section{Exercises}

\subsection*{Calculation}\pathtag{core}

\begin{exercise}
A promoter produces $r = 5$ mRNA per cell per generation. Each mRNA
is translated $30$ times. The cell divides every 40 minutes. The
protein is degraded only by dilution. Compute the steady-state
protein number per cell.

\paragraph{Solution sketch.}
Protein production per cell per generation: $5 \times 30 = 150$
proteins. Per-minute production: $150/40 = 3.75$ proteins/min.
Dilution rate: $\ln 2 / 40 = 0.01733\,\text{min}^{-1}$. At
steady state $r_{\text{prod}} = \mu N$, so $N = 3.75/0.01733 \approx
216$ proteins per cell. The simpler bookkeeping: dilution-limited
steady state equals total production over $1/\ln 2 \approx 1.44$
generations of accumulation, $150 \times 1.44 \approx 216$.
\end{exercise}

\begin{exercise}
For the toggle switch with $\alpha_1 = \alpha_2 = 5$ and Hill
exponents $m = n = 2$, find the three equilibrium points and
classify them (stable nodes or saddle points).

\paragraph{Solution sketch.}
At equilibrium $u^{*} = 5/(1+v^{*2})$ and $v^{*} = 5/(1+u^{*2})$.
By symmetry there is a solution on the diagonal $u^{*} = v^{*} = p$
satisfying $p(1+p^{2}) = 5$, so $p^{3} + p - 5 = 0$ with real
root $p \approx 1.516$. Off-diagonal solutions: substituting
$v^{*} = 5/(1+u^{*2})$ into the first equation and clearing,
$u^{*}(1 + 25/(1+u^{*2})^{2}) = 5$, i.e.\
$u^{*}\big((1+u^{*2})^{2} + 25\big) = 5(1+u^{*2})^{2}$. Numerical
solution gives the two off-diagonal equilibria
$(u^{*}, v^{*}) \approx (4.84, 0.198)$ and $(0.198, 4.84)$ by
symmetry. Linearisation: the Jacobian is
$J = \begin{pmatrix}-1 & -2v\alpha_1/(1+v^{2})^{2} \\ -2u\alpha_2/(1+u^{2})^{2} & -1\end{pmatrix}$.
At the diagonal point: $J_{12} = J_{21} = -2 \cdot 1.516 \cdot 5/(1+1.516^{2})^{2} \approx -1.00$;
eigenvalues $-1 \pm 1.00$, so one positive (unstable), one negative;
the diagonal point is a saddle. At
$(4.84, 0.198)$: $J_{12} = -2 \cdot 0.198 \cdot 5/(1.039)^{2} \approx -1.83$ and
$J_{21} = -2 \cdot 4.84 \cdot 5/(24.4)^{2} \approx -0.081$;
eigenvalues $-1 \pm \sqrt{0.148}$, both negative ($\approx -0.61,
-1.39$); the point is a stable node. By symmetry $(0.198, 4.84)$
is also a stable node. The system is bistable, as predicted by
$\alpha = 5 > \alpha_{c} \approx 2$ for $n = 2$.
\end{exercise}

\begin{exercise}
The Cas9 off-target frequency at a particular site is $10^{-4}$ per
edit. The on-target frequency is 0.5. A clinical protocol edits
$10^8$ cells. Compute the expected number of off-target edits and
the expected number of on-target edits.

\paragraph{Solution sketch.}
On-target: $0.5 \times 10^{8} = 5 \times 10^{7}$ edited cells.
Off-target: $10^{-4} \times 10^{8} = 10^{4}$ cells carry the
off-target edit. The ratio of on-target to off-target events is
$5000:1$. For a clinical product, the question is not the ratio but
the absolute frequency of off-target events in the infused
population: $10^{4}$ cells with an unwanted edit is large enough
that, if the off-target site is within a tumour-suppressor or
proto-oncogene, the long-term cancer risk is non-negligible.
Mitigations to acceptable rates require either a more specific
editing tool (high-fidelity Cas9, base editor) or a target whose
off-target sites are demonstrably benign by long-read or whole-genome
sequencing of the final product.
\end{exercise}

\begin{exercise}
A directed-evolution library has $10^6$ variants per round, each
with a probability $10^{-3}$ of having improved activity. Compute the
expected number of improved variants per round and the probability of
at least one improved variant.

\paragraph{Solution sketch.}
Expected improved variants per round: $10^{6} \times 10^{-3} =
10^{3}$. Probability of zero improved variants:
$(1 - 10^{-3})^{10^{6}} \approx \exp(-10^{6} \cdot 10^{-3}) =
\exp(-1000) \approx 10^{-434}$. Probability of at least one improved
variant: $1 - \exp(-1000) \approx 1$ to machine precision. The
selection saturates because the library is far larger than the
inverse hit rate. The practical bottleneck shifts to characterising
and ranking the thousand improved variants and choosing which to
carry to the next round; library design (focus mutagenesis on
positions of high impact) reduces the round size for a given
expected number of improved variants.
\end{exercise}

\begin{exercise}
Compute the dilution-limited protein degradation rate for a cell
dividing every 80 minutes. Compute the steady-state protein number
for a production rate of 100 proteins per generation.

\paragraph{Solution sketch.}
Dilution rate $\mu = \ln 2 / 80 = 0.00866\,\text{min}^{-1}$.
Production rate $r_{\text{prod}} = 100/80 = 1.25$ proteins/min.
Steady-state number $N = r_{\text{prod}}/\mu = 1.25/0.00866 \approx
144$ proteins per cell. Equivalently $100/\ln 2 \approx 144$ since
the steady state is $(1/\ln 2)$ times the per-generation
production. The slow-growth case (80-minute division) carries
roughly twice as many protein molecules per cell at the same
per-generation production as a fast-growth (40-minute division) case,
because the dilution rate is halved. The result is the engineering
intuition that strong promoters in slow-growing cells reach the
highest absolute protein numbers, but they also pay the largest
cellular-resource cost.
\end{exercise}

\subsection*{Derivation}\pathtag{standard}

\begin{exercise}
Derive the bistability condition for the toggle switch from
$\alpha_1, \alpha_2, m, n$. Identify the locus in parameter space
where bistability is lost (the saddle-node bifurcation).

\paragraph{Solution sketch.}
Substitute the second nullcline into the first to obtain a
fixed-point equation in one variable:
$u = \alpha_1/(1 + (\alpha_2/(1+u^{m}))^{n})$, i.e.\
$F(u) = u(1 + u^{m})^{n} + u \alpha_2^{n} - \alpha_1 (1 + u^{m})^{n} = 0$.
The number of real roots equals the number of fixed points. Three
real roots (the bistable case) require that $F(u)$ have two real
critical points, which by Rolle's theorem requires the derivative
$F'(u)$ to change sign twice. For the symmetric case
$\alpha_1 = \alpha_2 = \alpha$ and $m = n$, the analysis reduces by
substituting $u = v$ on the symmetric fixed point. The
saddle-node bifurcation locus is the set of $(\alpha_1, \alpha_2, m, n)$
where two of the three roots coalesce: $F(u_{*}) = 0$ and
$F'(u_{*}) = 0$ at the same point. For the symmetric case the
condition is $\alpha = \alpha_{c}(n)$, where
$\alpha_{c} = n/(n-1) \cdot ((n-1))^{1/n}$ for the saddle-node at
the symmetric point; explicitly $\alpha_{c}(2) = 2$ and
$\alpha_{c}(3) = 3/2 \cdot 2^{1/3} \approx 1.89$. Bistability
requires both $n > 1$ and $\alpha > \alpha_{c}$. Loss of cooperativity
($n \to 1$) eliminates bistability for any finite $\alpha$ because
the production rate then varies smoothly with the repressor
concentration and the system has a unique stable point.
\end{exercise}

\begin{exercise}
Derive the repressilator oscillation condition. Use linear stability
analysis at the symmetric equilibrium and identify the condition on
the Hill cooperativity and the production rate.

\paragraph{Solution sketch.}
Symmetric equilibrium $p_{1} = p_{2} = p_{3} = p^{*}$ satisfies
$p^{*} = \alpha_{0} + \alpha/(1 + p^{*\,n})$. The Jacobian of the
three-dimensional system at $p^{*}$ has block-circulant structure:
$J = -I + R$ where $R$ is the cyclic-shift matrix with diagonal
entry zero and one off-diagonal entry equal to
$-X = -n \alpha p^{*\,(n-1)}/(1 + p^{*\,n})^{2}$. Eigenvalues of
$R$: $-X \cdot \omega^{k}$ for $k = 0, 1, 2$ where
$\omega = e^{2\pi i/3}$. Eigenvalues of $J$: $-1 - X \omega^{k}$.
The two complex eigenvalues are $-1 - X(-1/2 \pm i\sqrt{3}/2) =
(-1 + X/2) \mp iX\sqrt{3}/2$. The Hopf bifurcation occurs when the
real part vanishes: $-1 + X/2 = 0$, i.e.\ $X = 2$. Substituting
$X$, the condition is $n \alpha p^{*\,(n-1)}/(1+p^{*\,n})^{2} = 2$.
For strong repression ($p^{*\,n} \gg 1$ so $p^{*} \approx \alpha_{0}$
and the right-hand denominator simplifies), the condition reduces
approximately to $n \alpha \gtrsim 2 \alpha_{0}^{2-n+1}$, requiring
both sharp cooperativity ($n > 2$, in practice) and high
production rate $\alpha$ relative to leakage. The original
Elowitz-Leibler repressilator with $n \approx 2$ and
$\alpha \approx 20 \cdot \alpha_{0}$ sat well inside the
oscillatory regime; the period scaled approximately as
$2 \pi / \mathrm{Im}(\lambda) = 4\pi/(X\sqrt{3})$ in
dimensionless units, multiplied by the protein decay time to
recover physical period.
\end{exercise}

\begin{exercise}
Derive the FBA linear program from the steady-state condition
$\mathbf{S}\mathbf{v} = \mathbf{0}$ and a biomass objective. Set up
the program for a toy network with three reactions and two
metabolites.

\paragraph{Solution sketch.}
The LP is
\begin{align*}
\text{maximise} & \quad c^{T} \mathbf{v} \\
\text{subject to} & \quad \mathbf{S} \mathbf{v} = 0, \\
& \quad v_{i}^{\min} \le v_{i} \le v_{i}^{\max}, \quad i = 1, \ldots, n,
\end{align*}
where $c$ is the objective vector (a unit weight on the biomass
flux for growth maximisation; a unit weight on the product flux
for product maximisation). Toy example: two metabolites $A$ (glucose,
supplied at rate $v_{1}$) and $B$ (target product). Three reactions:
$v_{1}$: glucose uptake, $v_{1} \in [0, 10]$ mmol gDW${}^{-1}$ h${}^{-1}$.
$v_{2}$: $A \to B$ at $v_{2} \in [0, \infty)$.
$v_{3}$: $B$ excretion at $v_{3} \in [0, \infty)$. Stoichiometric
matrix $\mathbf{S} = \begin{pmatrix}1 & -1 & 0 \\ 0 & 1 & -1\end{pmatrix}$.
Steady-state: $v_{1} = v_{2}$, $v_{2} = v_{3}$, so all three are
equal. Objective $\max v_{3}$ gives $v_{1} = v_{2} = v_{3} = 10$,
the maximum yield case. Adding a competing reaction $v_{4}$:
$A \to \text{biomass}$ at flux $v_{4}$ subject to a growth
requirement $v_{4} \ge v_{4}^{\min}$ produces the realistic
trade-off: maximum product flux is now
$v_{3} = v_{1} - v_{4}^{\min}$. The structure scales: full
genome-scale FBA models have 1000s of reactions and 100s of
metabolite balances, but the LP itself remains tractable in modern
solvers.
\end{exercise}

\begin{exercise}
Derive the noise amplification of a feedback genetic circuit: the
variance in protein concentration as a function of the production
and degradation rates. Use the Langevin approximation to the
chemical-master equation.

\paragraph{Solution sketch.}
For an unregulated gene with production rate $\beta$ and degradation
rate $\gamma$, the Langevin equation is
$dN/dt = \beta - \gamma N + \xi(t)$ with white-noise intensity
$\langle \xi(t) \xi(t') \rangle = (\beta + \gamma \langle N \rangle) \delta(t-t')$
(both production and degradation contribute the shot noise). At
steady state $\langle N \rangle = \beta/\gamma$ and the variance is
$\sigma^{2} = \langle N \rangle$, giving Fano factor $F = 1$
(Poisson). With negative feedback (production rate
$\beta(N) = \beta_{0}/(1 + N/K)$ for a simple repressor), linearise:
$\beta'(\langle N\rangle) = -\beta_{0}/K(1+\langle N\rangle/K)^{2}$.
The effective degradation rate is $\gamma_{\text{eff}} = \gamma - \beta'(\langle N\rangle) > \gamma$,
so the relaxation is faster and the variance falls to
$\sigma^{2} = (\beta + \gamma \langle N \rangle)/(2\gamma_{\text{eff}}) <
\langle N \rangle / 2$. Negative feedback divides the variance by
the open-loop gain $G = 1 + |\beta'(\langle N\rangle)| / \gamma$.
For a strong feedback ($G \sim 10$) the noise floor of an
otherwise Poisson process drops by a factor of 10, reaching
$F \approx 0.1$. This is Alon's ``noise filtering by negative
autoregulation,'' which is one reason naturally occurring
transcription factors so often regulate their own expression.
\end{exercise}

\subsection*{Estimation}\pathtag{core}

\begin{exercise}
Estimate the number of CRISPR edits required to disrupt a single
copy of the human genome (3.2 Gb) at every position, given a
20-nucleotide guide with PAM constraint.

\paragraph{Solution sketch.}
PAM constraint reduces the unique 20-mer target space by
$\approx 16$ (one accessible PAM per 16 bp on either strand).
Sites in the human genome: $\sim 3.2 \times 10^{9}/16 \approx 2 \times 10^{8}$
PAM-accessible positions. Each requires a distinct guide; saturation
of the genome thus requires $\sim 2 \times 10^{8}$ different guides.
At commercial pricing as of 2026 ($\sim \$15$ per chemically synthesised
sgRNA in 96-well-format batches, or $\sim \$0.005$ per sgRNA in
oligo-pool array synthesis), the array-pool cost is
$\sim \$10^{6}$ for the synthesis alone, plus library cloning,
delivery, and screening. Practical CRISPR screens (Sanjana 2014
GeCKO, Wang 2014 sgRNA pools) cover the protein-coding genome
($\sim 4 \times 10^{4}$ guides) rather than every PAM-accessible
position because saturation of intergenic regions is rarely worth
the cost. The order-of-magnitude answer: full-genome saturation is
on the boundary of practical, and the field has chosen targeted
libraries rather than saturation as the default.
\end{exercise}

\begin{exercise}
Estimate the time required to evolve a 100$\times$ improvement in
enzyme activity by directed evolution, assuming 2$\times$ per round,
a round time of two weeks, and library size sufficient for
2$\times$ improvement per round.

\paragraph{Solution sketch.}
Number of rounds: $\log_{2}(100) \approx 6.6$, so 7 rounds.
Wall-clock: $7 \times 2 = 14$ weeks. With a typical
post-characterisation gap of one to two weeks between rounds the
realistic timeline is closer to 20--25 weeks (5--6 months) per
project. PACE-equivalent timescale reduces this to two to four
weeks total for the same fold improvement, when the selection
can be coupled to phage propagation; the trade-off is loss of
visibility into the intermediate variants, which a discrete-round
campaign can characterise mid-stream.
\end{exercise}

\begin{exercise}
Estimate the cost of a clinical CRISPR therapy that requires
$10^{10}$ edited cells per patient, given a per-cell editing cost of
\$$10^{-7}$ in research-scale equipment and a regulatory and
manufacturing overhead of $1000\times$.

\paragraph{Solution sketch.}
Raw editing cost: $10^{10} \times 10^{-7} = \$1{,}000$ per patient.
With $1000\times$ overhead for cGMP manufacturing, sterility, lot
release, regulatory documentation, supply chain, and clinical
delivery: $\$1{,}000{,}000$ per patient. The 2024 list price of
Casgevy was \$2.2 million per patient
\cite{web:v6c09-fda-casgevy-2023}; the order-of-magnitude estimate
matches and the additional factor of two captures the
amortised R\&D recovery and the manufacturing capacity
constraint (Vertex and CRISPR Therapeutics' Casgevy is produced
in a small number of authorised treatment centres at limited
throughput). The implication for engineering planning: a
million-cell therapy costs roughly a million dollars per patient
under current manufacturing; a hundred-thousand-cell therapy might
reach \$100{,}000, a price compatible with broader payer coverage.
The engineering target for the next decade is to reduce per-cell
manufacturing cost by two orders of magnitude (closed-system
automated manufacturing, in vivo delivery that bypasses ex vivo
cell production) before CRISPR therapies become affordable at
population scale.
\end{exercise}

\subsection*{Simulation}\pathtag{mastery}

\begin{exercise}
Simulate the toggle switch by numerical integration. Sweep the
parameter $\alpha$ from 1 to 20 and identify the bistability range.
Plot the bifurcation diagram.

\paragraph{Solution sketch.}
Use `toggle\_switch\_dynamics.py` in the chapter's \texttt{code/}
directory; the script integrates the symmetric model
$du/d\tau = \alpha/(1+v^{n}) - u$ and analogous for $v$, sweeps
$\alpha \in [1, 20]$ at $n = 2$, and locates fixed points by
Newton iteration. Expected bifurcation diagram: a single stable
branch on the diagonal for $\alpha < \alpha_{c}(2) = 2$; a
pitchfork-like split into two off-diagonal stable branches plus the
diagonal saddle for $\alpha > 2$; the off-diagonal branches grow
as $u^{*} \approx \alpha - 1/\alpha^{n-1}$ for large $\alpha$.
Practical check: integrate from a perturbation off the diagonal at
$\alpha = 5$ and confirm convergence to one of the off-diagonal
stable points; integrate from exactly on the diagonal and confirm
the symmetry-preserving (saddle) trajectory.
\end{exercise}

\begin{exercise}
Simulate the repressilator with three repressors. Vary the
production rate and the Hill cooperativity. Report the period of
oscillation as a function of these parameters.

\paragraph{Solution sketch.}
Use `repressilator.py` in the \texttt{code/} directory; integrates the
three-protein ODE with parameter sweeps over
$\alpha \in [5, 200]$ and $n \in [1.5, 6]$. Detect oscillation by
running for 2000 dimensionless time units, discarding transient,
and applying Fourier analysis to the protein trajectory. Expected
results: no oscillation for $n \le 2.0$ regardless of $\alpha$;
oscillation onset at $n \approx 2$, $\alpha > 5$; period scales as
$\sim 2\pi/\sqrt{X-2}$ near onset, increasing as
$\alpha \to \alpha_{c}$ from above. Far from onset
($\alpha = 100$, $n = 4$), the period is approximately
$5\tau_{\text{prot}}$, where $\tau_{\text{prot}}$ is the protein
decay time. Compare with Elowitz-Leibler 2000 (Fig.~3) and
Potvin-Trottier et al.\ 2016 (Fig.~2c) for the
experimentally observed period range.
\end{exercise}

\begin{exercise}
Build a stochastic simulation of a synthetic feedback circuit using
the Gillespie algorithm. Compare the variance in protein
concentration to the deterministic prediction.

\paragraph{Solution sketch.}
Reaction set: $\emptyset \to N$ at rate $\beta(N)$; $N \to \emptyset$
at rate $\gamma N$. Implement the standard direct-method Gillespie:
sample exponential time to next event with rate
$a_{0} = \beta(N) + \gamma N$, choose event proportionally, update
$N$, advance time. Run to steady state, sample many independent
trajectories. Expected results for unregulated production
($\beta(N) = \beta_{0}$): mean $\langle N\rangle = \beta_{0}/\gamma$,
variance $\sigma^{2} = \langle N\rangle$, Fano $F = 1$ (Poisson).
For negatively autoregulated production
($\beta(N) = \beta_{0}/(1+N/K)$): mean
$\langle N\rangle < \beta_{0}/\gamma$, Fano $F < 1$, with reduction
factor of $\sim 1/(1+G)$ where $G$ is the open-loop gain at the
operating point. Match against the Langevin prediction from the noise-amplification
derivation exercise (within $\sim 10\%$ in the regime where the
linearisation is valid).
\end{exercise}

\subsection*{Design}\pathtag{standard}

\begin{exercise}
Design a genetic AND gate that produces fluorescence only when two
small molecules are present. Specify the regulatory architecture,
the parts, and the predicted dynamic range.

\paragraph{Solution sketch.}
Architecture: split-transcription factor. T7 RNA polymerase is split
into two complementary fragments T7-N and T7-C; each fragment is
expressed only in the presence of one inducer. T7-N expression is
driven by an arabinose-inducible promoter ($P_{\text{BAD}}$
controlled by AraC); T7-C expression is driven by an IPTG-inducible
promoter ($P_{\text{Lac}}$ controlled by LacI). Both fragments must
be present to reconstitute a functional T7 polymerase, which drives
expression of GFP from a T7 promoter. Parts (from the Registry or
addgene equivalents): pBAD-T7N-Lva (degradation-tagged for
responsiveness), pLac-T7C-Lva, T7p-sfGFP. Predicted dynamic range:
$\sim 100\times$ between fully induced and uninduced states, limited
by the leakage of $P_{\text{BAD}}$ (without arabinose, residual
T7-N expression at $\sim 1\%$ of induced level) and of
$P_{\text{Lac}}$ ($\sim 0.5\%$ without IPTG, $\sim 0.1\%$ with
glucose catabolite repression). Time to steady state: $\sim 3$
generations after induction ($\sim 2$ hours in fast-growing
\emph{E.\ coli}); reset time after induce removal: $\sim 4$
generations due to T7 fragment dilution. Failure modes: T7 polymerase
toxicity at high expression (cured by Lva-tag degradation);
inducer cross-talk (cured by careful titration of inducer
concentrations to operating range).
\end{exercise}

\begin{exercise}
Design a metabolic-engineering route from glucose to lactate in
\emph{E.\ coli}, identifying the principal pathway, the
gene-knockout targets to prevent off-pathway loss, and the predicted
maximum yield.

\paragraph{Solution sketch.}
Pathway: glycolysis from glucose to pyruvate (10 reactions, net
$\text{glucose} + 2\,\text{NAD}^{+} + 2\,\text{ADP} + 2\,P_{i}
\to 2\,\text{pyruvate} + 2\,\text{NADH} + 2\,\text{ATP}$), then
lactate dehydrogenase (LDH, gene \emph{ldhA} in \emph{E.\ coli})
converting pyruvate to lactate with NAD${}^{+}$ regeneration:
$\text{pyruvate} + \text{NADH} \to \text{lactate} + \text{NAD}^{+}$.
Maximum theoretical yield: 2 mol lactate per mol glucose, i.e.\
$2 \cdot 90/180 = 1.0$ g/g. Knockouts to prevent off-pathway flux:
\emph{pta} and \emph{ackA} (acetate production from acetyl-CoA);
\emph{adhE} (ethanol production); \emph{frdA} (succinate production
under anaerobic conditions); \emph{poxB} (pyruvate oxidase).
Over-expression: \emph{ldhA} from a strong promoter. The Zhou et al.\
(2003) and Mazumdar et al.\ (2010) engineered \emph{E.\ coli} strains
reach $\sim 0.95$ g/g yield in anaerobic batch fermentation,
within 5\% of the theoretical maximum. Process: anaerobic shake
flask or stirred tank; pH control with $\text{Ca(OH)}_{2}$ to
neutralise lactic acid; harvest before stationary phase to avoid
lactate re-utilisation. Output measurement: HPLC for lactate
quantification, OD$_{600}$ for biomass.
\end{exercise}

\begin{exercise}
Design a CRISPR base-editing strategy to convert a single
nucleotide in a clinical target. Specify the editor (cytidine or
adenosine), the guide RNA, the delivery method, and the predicted
on-target and off-target frequencies.

\paragraph{Solution sketch.}
Example target: correct the sickle-cell mutation (HBB Glu6Val,
$\text{GAG} \to \text{GTG}$) by reverting $T \to A$, i.e.\ an
A-to-G or T-to-C base conversion. Cytidine base editors (CBE)
convert $C \to T$; adenine base editors (ABE) convert $A \to G$.
The reversion of T (template strand) to A is performed by an ABE
acting on the complementary strand (where the mutated base reads
as A in the template strand and the conversion is A-to-G read on
the antisense). Editor: ABE8e (Richter et al.\ 2020), which has
faster kinetics and higher activity than the original ABE7.10.
Guide RNA: 20 nt protospacer placing the target A in positions
4--8 (the editor's narrow editing window) relative to the PAM. PAM
constraint: SpCas9 NGG; for HBB the available PAMs constrain the
guide placement to one of two viable sgRNAs, both verified for
single-A specificity within the window. Delivery: ex vivo electroporation
of editor mRNA and sgRNA into autologous CD34+ hematopoietic stem
cells, followed by re-infusion (autologous gene-therapy paradigm
matching Casgevy). On-target editing efficiency: $\sim 70$--90\%
A-to-G at the target position in HSPCs in single-cell readout.
Off-target editing: bystander A conversion within editing window
on the same guide ($\sim 5$--20\% depending on the adjacent A);
genome-wide off-target by ABE deamination of A bases outside the
guide context, estimated at $\sim 0.01$\% per A per cell genome-wide,
which gives $\sim 3 \times 10^{5}$ stochastic deaminations per cell
across the $3 \times 10^{9}$-bp genome. The number is large; the
clinical question is what fraction of those off-target deaminations
land in coding regions or splice sites; the answer is $\sim 1$--2\%,
so $\sim 3000$ functional bystander edits per cell, of which the
distribution of consequence is dominated by intronic positions.
The current frontier is engineering tighter ABE variants
(ABE9, evolved by PACE) and verifying the off-target profile in
each clinical target.
\end{exercise}

\subsection*{Judgement}\pathtag{core}

\begin{exercise}
A laboratory proposes to construct an \emph{E.\ coli} strain that
synthesises a peptide hormone for direct human dietary use. State the
regulatory framework that applies, the safety studies required, and
the engineering specifications you would impose on the strain
beyond those required for laboratory production of the hormone.

\paragraph{Discussion.}
Regulatory frameworks: in the US, the FDA Center for Drug Evaluation
and Research (CDER) for the peptide as a drug, plus the FDA Center
for Food Safety and Applied Nutrition (CFSAN) and USDA APHIS for
any environmental release of a live engineered organism; in the EU,
EMA and EFSA jointly, with NCB (National Competent Body) involvement
in each member state. Safety studies: full pharmacology (peptide
identity, purity, dose-response, pharmacokinetics) plus toxicology
(acute, sub-chronic, reproductive); allergenicity assessment of the
peptide if novel; absence of bacterial residues (endotoxin, LPS,
host-cell protein), absence of the bacterial DNA in the product
(or characterisation of horizontal-transfer risk in the gut);
microbiome compatibility studies if the strain is intended to
survive in the gut. Engineering specifications beyond laboratory
production: kill switch triggered by ambient (non-laboratory)
conditions (e.g., absence of an inducer present in the laboratory
medium but not in the gut); auxotrophy for a non-natural amino acid
or a synthetic vitamin not present in food; genome reduction to
remove pathogenicity factors, antibiotic-resistance markers, and
mobile genetic elements; verification that the engineered strain
cannot transfer the peptide-expression cassette to wild-type gut
bacteria; biocontainment specifications meeting Lawrence Berkeley
Lab BSL-1-CAT or stricter; release-monitoring plan if any escape
is possible. The strain is a living biological product, not a
bulk chemical; regulatory and engineering specifications must
match the higher hazard of a self-replicating organism.
\end{exercise}

\begin{exercise}
A gain-of-function proposal seeks to identify mutations in an
animal coronavirus that confer human-cell receptor binding. The
proposal includes BSL-3 containment, a claimed vaccine-development
benefit, and no immediate therapeutic product. State the engineering
questions you would require answered before the work proceeds: failure
mode, containment layer, operator exposure, environmental release,
information hazard, and stopping rule.

\pa

\begin{exercise}
A proposed environmental biosensor strain is designed to fluoresce in
the presence of arsenic in groundwater. State the containment design,
the retrieval plan, the kill-switch verification test, and the
measurement that would convince you the field deployment is failing.
\end{exercise}

\begin{exercise}
A metabolic-engineering team reports a strain that produces a target
chemical at high titer in shake flasks but loses 70\% of the titer in
a stirred bioreactor. State three hypotheses that connect the genetic
design to the process environment, and state one measurement for each.
\end{exercise}

\begin{exercise}
A company proposes to market a genetically modified probiotic that
secretes a therapeutic peptide in the gut. State the safety evidence,
genetic-stability evidence, dosing-control evidence, and post-release
monitoring you would require before treating it as an engineered
product rather than a laboratory demonstration.
\end{exercise}

\begin{exercise}
A CRISPR-edited crop has a desired yield trait and an unexpected
change in flowering time. State how you would distinguish an off-target
edit, a linked on-target effect, and an environmental interaction.
\end{exercise}

\begin{exercise}
A synthetic-biology project claims that standard parts make the design
portable across \emph{E.\ coli}, yeast, and a mammalian cell line.
State which parts of the claim are biologically implausible without
new characterisation data, and name the data you would request.
\end{exercise}
